% ch-nonparam --- nonparametric intensities: boosted trees, forests, small networks

\section{Scope and placement}\label{chnp:sec:scope}

This chapter keeps the estimation principle of
\cref{ch:poisson,ch:hawkes} and relaxes the functional form. The
intensity is still fitted by maximum likelihood --- the objective is
still the point-process log-likelihood of \cref{thm:ch03:jacod} --- but
the intensity is no longer log-linear in the features. The model is
\begin{equation}\label{eq:chnp:model}
\lam_i(t) \;=\; g\big(f(w_i(t))\big),
\end{equation}
where $f$ is a flexible regression function --- a gradient-boosted tree
ensemble, a random forest, or a small neural network --- and $g$ is a
fixed positivity link, either $g(z) = e^{z}$ or the softplus
$g(z) = \log(1 + e^{z})$ discussed for the parametric case in
\cref{ch04:sec:link}. The feature vector $w_i(t)$ has two blocks:
\begin{enumerate}[label=(\roman*)]
\item the point-in-time covariates $x_i(t)$ of \cref{prot:ch02:pit}
(macro series, sector dummies, calendar terms, firm attributes); and
\item the \emph{engineered history features} of
\cref{ch05:sec:inhibition}: the time since the firm's last deal, decayed
event counts $H_{\decay}(t)$ at several fixed half-lives
(\cref{eq:ch05:decayedcount}), and recent-deal counts in the firm's
sector.
\end{enumerate}
Nothing in block (ii) is new. These are the same excitation and
inhibition covariates that route (a) of \cref{ch:hawkes} feeds into a
log-linear intensity. What changes is the map from features to
log-intensity: $f$ may bend, saturate, and interact them freely.

\begin{assumption}[Trusted, uncensored deal times]\label{asm:chnp:trusted}
Event times are correct at the resolution of the analysis bins, and no
event is interval-censored beyond that resolution. Both the bin
assignment of each event and the engineered history features of
\cref{ch05:sec:inhibition} consume exact times; neither survives date
corruption.
\end{assumption}

\begin{warning}[This regime has no defense against untrusted
dates]\label{warn:chnp:dates}
The parametric chapters carry an escape route: when dates are
unreliable, the interval-censored likelihoods of \cref{ch:censoring}
replace the exact likelihood. That route depends on closed-form
compensators, which \cref{eq:chnp:model} does not have. If
\cref{asm:chnp:trusted} fails --- vendor-imputed dates, month-rounded
dates, bulk backfills --- do not fit this chapter's models. Fall back to
\cref{ch:censoring}, or coarsen the bins to the trusted resolution and
drop every feature built at a finer one.
\end{warning}

These models drop the shape assumptions of the parametric chapters and
give up most of what those assumptions buy. Their role is the
\emph{challenger} (\cref{chnp:sec:evaluation}).

\section{The reduction to Poisson regression}\label{chnp:sec:reduction}

Machine-learning software minimizes losses over rows; it does not
maximize point-process likelihoods. One reduction closes the gap: on a
fine enough grid, the point-process likelihood \emph{is} a Poisson
regression likelihood over (unit, bin) rows with exposure offsets. This
is the point-process GLM framing of \citet{truccolo2005};
\cref{thm:ch04:glmequiv} is its log-linear special case, and the version
we need allows an arbitrary $f$.

Partition $(0,\Tend]$ into bins $B_k = (a_{k-1}, a_k]$, $k = 1,\dots,K$,
with widths at most $\delta$. For each unit $i$ and bin $k$ define the
count, the exposure, and the bin feature
\begin{equation}\label{eq:chnp:rows}
y_{ik} \;=\; \#\{m : t_m \in B_k,\ i_m = i\},
\qquad
e_{ik} \;=\; \int_{B_k} \atrisk_i(u) \dd u,
\qquad
w_{ik} \;=\; w_i(a_{k-1}^{+}) .
\end{equation}

\begin{assumption}[Piecewise-constant features]\label{asm:chnp:pwconst}
$w_i(\cdot)$ is constant on every bin $B_k$: covariates are LOCF series
aligned to the grid, and history features are frozen at each bin's left
endpoint. Under \cref{eq:chnp:model} the intensity is then a step
function with values $\lam_{ik} = g(f(w_{ik}))$.
\end{assumption}

\begin{proposition}[Binned Poisson reduction]\label{prop:chnp:binned}
Under \cref{asm:chnp:pwconst}:
\begin{enumerate}[label=(\roman*)]
\item The point-process log-likelihood of \cref{thm:ch03:jacod} equals
\begin{equation}\label{eq:chnp:binnedloglik}
\loglik(f)
 \;=\;
 \sum_{i,k} \Big( y_{ik} \log \lam_{ik} \;-\; e_{ik}\, \lam_{ik} \Big),
\end{equation}
exactly, and \cref{eq:chnp:binnedloglik} differs from the log-likelihood
of independent counts $y_{ik} \sim \mathrm{Poisson}(e_{ik}\lam_{ik})$ by
an additive constant free of $f$.
\item Consequently, any learner that minimizes the Poisson deviance over
the rows \cref{eq:chnp:rows}, with $\log e_{ik}$ as an offset (or
exposure weight, in the equivalent rate parameterization), performs
maximum likelihood for the point process within the class of intensities
measurable with respect to the binned features.
\item If the true feature path updates continuously between grid points
(history features decay as in \cref{eq:ch05:decayedcount}) and is frozen
at bin left endpoints, then on every bin containing at most one event
the log-likelihood error of the frozen model relative to the
continuous-feature model is $O(\delta)$, uniformly over intensity paths
bounded away from $0$ and $\infty$.
\end{enumerate}
\end{proposition}

\begin{proof}
(i) The intensity is constant on each bin, so the compensator integral
splits into $\sum_{i,k} \lam_{ik} \int_{B_k} \atrisk_i =
\sum_{i,k} \lam_{ik} e_{ik}$. Each event of unit $i$ in bin $k$
contributes $\log\lam_i(t_m) = \log\lam_{ik}$ whatever its position in
the bin, so the event term groups into $\sum_{i,k} y_{ik}\log\lam_{ik}$.
The Poisson count log-likelihood is
$\sum_{i,k} [\, y_{ik}\log(e_{ik}\lam_{ik}) - e_{ik}\lam_{ik}
- \log y_{ik}! \,]$, which equals \cref{eq:chnp:binnedloglik} plus
$\sum_{i,k} [\, y_{ik}\log e_{ik} - \log y_{ik}! \,]$, a statistic of
the data alone.
(ii) The exposure-weighted Poisson deviance is, up to sign and a
data-only constant, the Poisson count log-likelihood of (i).
(iii) Between events the features move only by decay, so on $B_k$ the
continuous intensity satisfies
$\sup_{u \in B_k} \abs{\lam_i(u) - \lam_{ik}} = O(\delta)$ whenever
$g \circ f$ is Lipschitz on the (bounded) range of the features; the
compensator error on the bin is therefore at most $e_{ik} \cdot
O(\delta) = O(\delta)$. For the event term, with at most one event in
the bin, the feature seen by that event at $t_m^-$ involves no other
within-bin event, so the frozen value differs from the true value only
by within-bin decay drift, again $O(\delta)$; since the intensity is
bounded away from zero, $\log$ is Lipschitz on its range and the
log-intensity error is $O(\delta)$. A second event in the same bin would
break this: the first event's jump in the history features would be
missed. The bin width is therefore chosen below the shortest engineered
half-life, which also makes two-event bins rare.
\end{proof}

\begin{remark}[One table, three model families]\label{rem:chnp:onetable}
\Cref{prop:chnp:binned} is the whole interface. Build the (unit, bin)
table once and every model in this chapter trains on it, alongside the
log-linear GLM of \cref{ch:poisson} as the reference. All report the
same held-out point-process NLL, so the comparisons of
\cref{chnp:sec:evaluation} are exact.
\end{remark}

\section{Boosted trees}\label{chnp:sec:trees}

Gradient boosting fits $f$ as a sum of regression trees, each tree
fitted to the gradient of the loss at the current fit
\citep{friedman2001}. XGBoost implements the Poisson case directly
\citep{chenguestrin2016}: objective \code{count:poisson}, log link, and
\code{base\_margin} set to $\log e_{ik}$ so that the exposure enters as
a fixed offset rather than a feature.

The mechanics are transparent. Write $u_r = f(w_r)$ for the raw score of
row $r = (i,k)$, $o_r = \log e_r$ for the offset, and
$m_r = \exp(u_r + o_r) = e_r\, e^{u_r}$ for the predicted mean count.
The per-row negative log-likelihood, gradient, and Hessian in the score
are
\begin{equation}\label{eq:chnp:gradhess}
\ell_r(u_r) = m_r - y_r\,(u_r + o_r) + \log y_r! ,
\qquad
\frac{\partial \ell_r}{\partial u_r} = m_r - y_r ,
\qquad
\frac{\partial^2 \ell_r}{\partial u_r^2} = m_r .
\end{equation}
The gradient is ``predicted minus observed'' and the Hessian is the
predicted mean --- the same canonical-link structure as the GLM score
\cref{eq:ch04:score}, evaluated tree by tree. Two controls keep the
ensemble economically sensible and temporally honest.

\paragraph{Monotone constraints.}
Trees will fit a wiggle in any direction the training noise suggests.
Where economics fixes a direction, we impose it:
\code{monotone\_constraints} restricts every split so that $f$ is
monotone in the flagged feature. Candidates: intensity nonincreasing in
the policy rate for late-stage deal flow; nondecreasing in the sector's
recent-deal count. Features whose true shape we expect to be
non-monotone --- time since the firm's last deal, which should show
inhibition then recovery --- stay unconstrained; discovering that shape
is one reason to run this model. Constraints double as regularization
along directions where the sign is already known.

\paragraph{Early stopping on a temporal window.}
The number of boosting rounds is the main capacity knob. It is chosen by
early stopping on a validation window strictly later in time than the
training rows, matching \cref{ch07:sec:temporal}. Random row splits are
forbidden: rows from the same week share macro covariates and sector
history features, so a random split leaks the validation period into
training and overstates capacity.

\begin{protocol}[Boosted-tree intensity fit]\label{prot:chnp:xgb}
(a) Build the (unit, bin) table of \cref{eq:chnp:rows} with bin width at
most the shortest engineered half-life. (b) Fit XGBoost with objective
\code{count:poisson}, \code{base\_margin} $= \log e_{ik}$, monotone
constraints from the signed-feature registry, and early stopping on the
last temporal block of the training window. (c) Report held-out
point-process NLL via \cref{eq:chnp:binnedloglik}, PIT calibration, and
the time-rescaling diagnostics of \cref{chnp:sec:ledger}. (d) Report
uncertainty by the firm-cluster bootstrap of \cref{chnp:sec:ledger}.
\end{protocol}

Survival-flavored variants exist: accelerated-failure-time objectives in
XGBoost handle censored durations directly \citep{barnwal2022}, and
gradient-boosted Cox models are standard \citep{polsterl2020}. We do not
use them: the recurrent-event, time-varying-covariate structure of deal
flow maps more cleanly onto \cref{prop:chnp:binned}, which keeps the
likelihood identical to the rest of the document.

\section{Random forests}\label{chnp:sec:forests}

A random forest fits many deep trees on bootstrap resamples and averages
them. Two routes adapt forests to the (unit, bin) table. The clean route
is Poisson-deviance splitting: grow each tree with the Poisson deviance
as split criterion, weight rows by exposure $e_{ik}$, and predict the
rate $y/e$ within each leaf. The pragmatic route, for software without a
Poisson criterion, is exposure-weighted least squares on empirical log
rates $\log\{(y_{ik} + c)/e_{ik}\}$ with a small stabilizing constant
$c$; that is a quasi-likelihood approximation, not maximum likelihood,
and must be flagged as such in any NLL comparison.

An honest account of what averaging does. A forest prediction is an
average of training-set leaf values, so the fit never leaves the range
of what it saw: forests do not extrapolate. When a macro covariate exits
its training range, the response goes flat. That is a defect (a genuine
trend continues and the forest will not follow it) and a defense (the
log-linear GLM extrapolates \emph{exponentially} in the same situation;
\cref{warn:ch04:separation}), and flat is often the safer error. The
forest is therefore the stability-oriented challenger, the boosted
ensemble the capacity-oriented one, and disagreement between them
outside the training support is itself a diagnostic.

Out-of-bag evaluation scores each row with the trees that did not see
it, an almost-free deviance estimate. It is a screening tool only: OOB
rows are interleaved in time with training rows, so OOB deviance
inherits the temporal leakage that \cref{ch07:sec:temporal} exists to
prevent. Cross-family model selection still runs on the temporal splits;
OOB tunes forest internals (depth, feature subsampling) before that
comparison.

\section{Small neural networks}\label{chnp:sec:nets}

Neural networks enter twice, once cheaply and once exactly.

\paragraph{Route (i): the binned Poisson route.}
A small multilayer perceptron replaces the tree ensemble in
\cref{prop:chnp:binned}: input $w_{ik}$, a few narrow hidden layers,
and a softplus output head
\begin{equation}\label{eq:chnp:softplus}
\lam_{ik} \;=\; \log\big(1 + e^{z_{ik}}\big),
\qquad z_{ik} = \text{network}(w_{ik}),
\end{equation}
trained by minimizing the same exposure-weighted Poisson deviance. The
softplus head keeps the intensity positive without the exponential
amplification that \cref{warn:ch05:mechB} warns about. Everything from
the tree discussion transfers: temporal early stopping, the shared
table, the same held-out NLL. The network adds smooth interpolation and
costs optimization noise.

\paragraph{Route (ii): the exact continuous-time likelihood.}
Nothing forces binning. The exact objective is
\begin{equation}\label{eq:chnp:exactll}
\loglik(f)
 \;=\;
 \sum_{m} \log g\big(f(w_{i_m}(t_m^-))\big)
 \;-\;
 \sum_{i} \int_0^{\Tend} \atrisk_i(u)\,
 g\big(f(w_i(u))\big) \dd u ,
\end{equation}
and it is computable whenever the integral is. Between events the
engineered history features decay exponentially
(\cref{eq:ch05:decayedcount}) and the covariates are piecewise constant,
so $w_i(u)$ follows a known closed-form curve on each inter-event
interval. For log-linear $f$ the integral then has closed pieces --- the
parametric machinery of \cref{ch:hawkes} --- but for a generic network
$g(f(\cdot))$ along a known curve has no closed integral, and each
inter-event piece is evaluated by numerical quadrature (a few points per
piece suffice when the piece is shorter than the fastest half-life).
Automatic differentiation passes through the quadrature, so training is
standard.

Route (ii) is the door to neural temporal point processes, which
replace the engineered features themselves with a learned history
embedding: recurrent embeddings with an exponential-form intensity
\citep{du2016}, continuous-time recurrent intensities
\citep{meieisner2017}, networks for the compensator whose derivative
gives the intensity exactly \citep{omi2019}, and intensity-free models
of the inter-event distribution \citep{shchur2020}. These are the fully
flexible end of the axis this chapter walks. We do not start there, for
three plain reasons. Data volume: a typical firm contributes a handful
of deals, and learned embeddings need long streams or aggressive pooling
that reintroduces, in opaque form, the heterogeneity questions of
\cref{ch05:sec:frailty}. Interpretability: the engineered features are
the economic quantities we want shape estimates for --- cooldown,
momentum, sector spillover --- and an embedding answers none of them
directly. Calibration risk: flexible intensity models fail by confident
misallocation of intensity, and diagnosing that inside an embedding is
much harder than inspecting a feature's partial response. Neural TPPs
are the recorded escalation if the challenger wins and its wins trace to
history representation rather than covariate nonlinearity.

\section{What the regime keeps and what it loses}\label{chnp:sec:ledger}

\Cref{tab:chnp:ledger} states the trade explicitly.

\begin{table}[htbp]
\centering
\small
\begin{tabular}{p{0.24\textwidth} p{0.12\textwidth} p{0.52\textwidth}}
\toprule
\textbf{Property} & \textbf{Status} & \textbf{Consequence or replacement} \\
\midrule
Likelihood objective & kept & \cref{prop:chnp:binned}: the same
point-process NLL as \cref{ch:poisson,ch:hawkes}, so held-out
comparisons across regimes are exact \\
\addlinespace
Temporal evaluation & kept & the protocol of \cref{ch07:sec:temporal}
applies unchanged; early stopping and tuning obey it \\
\addlinespace
Time-rescaling diagnostics & kept & \cref{thm:ch03:rescaling} holds for
\emph{any} fitted intensity; see below \\
\addlinespace
Simulation & kept & thinning \citep{lewisshedler1979} with
piecewise-constant fitted rates; bound refresh at bin edges \\
\addlinespace
Branching interpretation & lost & no $\branch$, no endogenous share, no
parent attribution (\cref{thm:ch05:cluster} does not apply) \\
\addlinespace
Closed-form compensators & lost & compensators are bin sums or
quadrature; the interval-censored route of \cref{ch:censoring} is
unavailable \\
\addlinespace
Stability theory & lost & no spectral-radius certificate; only bounded
features and simulation monitoring (\cref{warn:ch05:mechB} discipline) \\
\addlinespace
Sandwich standard errors & lost & no low-dimensional score; use the
firm-cluster bootstrap \citep{efron1994} \\
\addlinespace
Frailty-vs-contagion reasoning & weakened & see
\cref{warn:chnp:contagion} and \cref{ch05:sec:frailty} \\
\bottomrule
\end{tabular}
\caption{What the nonparametric-intensity regime keeps from the
parametric chapters and what it gives up.}
\label{tab:chnp:ledger}
\end{table}

Two rows deserve expansion. \emph{Time rescaling survives everything.}
The time-rescaling theorem (\cref{ch03:sec:rescaling}) makes no
parametric assumption: if the fitted intensity is right, the rescaled
gaps are unit exponential, whether the intensity came from a GLM, a
Hawkes kernel, or a thousand trees \citep{brown2002,ogata1988}. The
fitted compensator here is a running sum of $\lam_{ik} e_{ik}$ over
bins, evaluated with frozen parameters on held-out windows as in
\cref{ch:gof}; this diagnostic holds the flexible models to the same
standard as the parametric ones. \emph{Uncertainty comes from the
bootstrap.} The sandwich variance of \cref{ch:foundations} needs a score
and a Hessian in a fixed finite-dimensional parameter, which a tree
ensemble does not have. The replacement is the cluster bootstrap
\citep{efron1994}: resample \emph{firms} with replacement, keeping each
firm's whole event history and exposure intact so within-firm dependence
is never broken; refit; read uncertainty off the resampled predictions.
Resampling rows would pretend the (unit, bin) cells are independent,
which they are not.

\section{Hawkes-flavored behavior without a Hawkes
model}\label{chnp:sec:imitation}

In this regime, excitation and inhibition enter through one door only:
the engineered history features. A positive learned response to the
decayed own-sector count looks like self-excitation; a depressed
response at small time-since-last-deal looks like inhibition. The
flexible $f$ can shape these responses --- thresholds, saturation,
interactions with the macro state --- which the log-linear form cannot.
That freedom is the point of the chapter, and it is also where the
causal reading breaks.

\begin{warning}[Learned excitation is not
contagion]\label{warn:chnp:contagion}
A latent firm-quality or sector-attractiveness factor that the
covariates miss makes past deal counts a proxy for the missing factor.
A flexible $f$ finds and uses that proxy more effectively than a linear
model: it can carve the feature space into regions where the proxy works
best and imitate self-excitation with no contagion present. This is the
confound of \cref{ch05:sec:frailty} with fewer restraints. The
parametric chapter could at least bound the damage --- branching
estimates are upper bounds, parent-attribution timescales provide a
probe --- and none of that structure exists here: no branching ratio, no
genealogy, no timescale decomposition. Any ``deals cause deals'' claim
from this regime is therefore strictly weaker than the corresponding
claim from the parametric model. The reporting rule: nonparametric fits
make \emph{predictive} statements about history features, never
mechanism statements; if the mechanism question matters, the parametric
layer answers it with the diagnostics of \cref{ch:hawkes}.
\end{warning}

\section{Evaluation and the challenger role}\label{chnp:sec:evaluation}

Everything runs inside the harness of \cref{ch07:sec:constitution}: the
same temporal splits, the same held-out point-process NLL (identical by
\cref{prop:chnp:binned}), the same calibration reports --- PIT
histograms for counts and proper scoring rules for predictive
distributions \citep{gneitingraftery2007}. Two comparisons are
designated: the log-linear GLM of \cref{ch:poisson} \emph{on the same
feature table}, so any gap is attributable to functional form rather
than information; and the parametric Hawkes layer of \cref{ch:hawkes},
so the value of genuine event-time structure is measured against
flexible use of engineered summaries.

The logic of the comparison is asymmetric. The nonparametric model is
not a candidate for the production intensity; it is the
\emph{challenger}, and its job is to expose structure the log-linear
form misses. If it beats the log-linear GLM materially on held-out NLL
and calibration, the conclusion is not ``deploy the trees''; it is that
the log-linear intensity is missing an interaction or a nonlinearity,
and the follow-up is to find it: partial-dependence and interaction
diagnostics on the fitted ensemble point at candidate terms, which are
added to the GLM explicitly and re-tested. If the challenger cannot beat
the GLM, the log-linear form is adequate at the current data size and
the parametric stack keeps its advantages.

\begin{decision}{D-NP.1 --- the nonparametric intensity is the
designated challenger}
The models of this chapter are fitted in every evaluation campaign, on
the shared (unit, bin) table, under the harness of
\cref{ch07:sec:constitution}, next to the log-linear GLM and the
parametric Hawkes layer. Their role is challenger, not candidate: a
material win over the log-linear intensity on held-out NLL and
calibration is treated as evidence of missing structure, triggering a
search (interaction and nonlinearity extraction from the ensemble, then
explicit GLM terms and a re-run), not a deployment. Rationale: the
challenger keeps the likelihood and the evaluation protocol
(\cref{prop:chnp:binned}), so its wins isolate functional form; the
interpretability, stability certificates, and inference that production
requires live in the parametric stack (\cref{tab:chnp:ledger}).
Reversal trigger: if the challenger's material win persists across two
consecutive campaigns \emph{after} the extracted terms are added to the
GLM, the deficit is structural, and promotion of a constrained tree or
network intensity to production candidate is put on the table, carrying
this chapter's bootstrap-inference and monitoring obligations with it.
\end{decision}

\section{Defect declaration}\label{chnp:sec:defects}

Per decision D2.1 (\cref{ch:data}), against the ledger of
\cref{tab:ch02:defects}.

\emph{(C1) Right censoring: corrected exactly.} The exposure $e_{ik}$
integrates the at-risk indicator to the end of observation, and
\cref{prop:chnp:binned} carries the survival factor through the bin
sums, as in \cref{ch:poisson}. Open spells must be kept; the zero-count
rows are the censored exposure.

\emph{(C2) Left truncation: corrected mechanically.} Exposure starts at
entry. The estimand remains conditional on entry into observation;
\cref{ch:universe} owns the unconditional question.

\emph{(C3) Unlabeled exits: vulnerable, downward.} Dead firms kept at
risk contribute exposure and no events, deflating fitted intensities.
The flexible $f$ adds a twist: it can learn that long-quiet firms have
low rates and partially \emph{mask} the defect in-sample without fixing
it --- the phantom exposure is still wrong. Treatments stay where they
were: exposure windows and mover--stayer structure
(\cref{ch:poisson,ch:buckets}), ranking-layer mitigations
(\cref{ch:ranking}).

\emph{(C4) Untrusted times: vulnerable by assumption.} This regime
consumes exact times twice, in the bin assignment and in the engineered
history features, and has no interval-censored fallback
(\cref{warn:chnp:dates}). \Cref{asm:chnp:trusted} is a precondition, not
a robustness claim; when it fails, \cref{ch:censoring} is the route.

\emph{(C5) Stale covariates: vulnerable, with a tree-specific hazard.}
Stale features attenuate learned responses as in \cref{ch:stale}. Trees
add a hazard the GLM lacks: staleness \emph{patterns} (reporting gaps,
vendor refresh cycles) are learnable signals, and a tree will exploit
them if they correlate with outcomes. The guards are the point-in-time
protocol \cref{prot:ch02:pit} and explicit staleness-age features, so
the dependence is visible rather than implicit.

\emph{(C6) Out-of-universe deals: sector features complete, firm
intensities thinned.} Sector-level history features are built from the
complete deal stream and remain valid. Firm-level intensities miss
unseen competitors exactly as in the parametric chapters;
\cref{ch:universe} owns the correction.

\begin{codehooks}
\emph{To be written.} \code{hawkes\_calibration/nonparam.py}: (1) the
(unit, bin) table builder of \cref{eq:chnp:rows}, shared with the GLM
fitters so that \cref{prop:chnp:binned}(ii) holds by construction; (2)
the XGBoost wrapper of \cref{prot:chnp:xgb} --- objective
\code{count:poisson}, \code{base\_margin} $= \log e_{ik}$, a
\code{monotone\_constraints} map read from the signed-feature registry,
temporal early stopping; (3) a forest wrapper with Poisson-deviance
splitting, exposure weights, and OOB reporting (screening only); (4) a
small-network trainer with the softplus head of \cref{eq:chnp:softplus}
and exposure-weighted Poisson loss; (5) a time-rescaling adapter that
consumes any \code{predict\_rate(unit, bin)} function and emits the
diagnostics of \cref{ch:gof}; (6) the firm-cluster bootstrap driver.
The \code{xgboost} objective configuration for (2) is a checked-in
parameter file, so challenger runs are reproducible across campaigns.
\end{codehooks}
